\section{Related Work}

The cost of security mechanisms has been widely studied, particularly in the context of network protocols. 
A key reference is \cite{naylor2014cost}, where the authors evaluated the cost of the "S" in HTTPS, demonstrating that adding security through TLS introduces additional data overhead and computational cost. This work provides a foundation for analyzing security not only as a functional property, but also as a measurable source of resource consumption.

In the context of Cyber-Physical Systems (CPS), Ivki\'c et al. \cite{ref01, ref02, ref03} propose a series of security cost modeling frameworks that highlight the importance of evaluating security mechanisms alongside their resource and operational overhead. First, Ivki\'c et al. \cite{ref01} introduced the concept of modeling security costs in terms of resource consumption, showing that such mechanisms can significantly affect constrained environments. Next, they extended this by proposing a runtime framework \cite{ref02}, demonstrating that security overhead varies dynamically with system state and workload. Finally, they presented a more comprehensive model integrating computational, communication, and operational costs \cite{ref03}, emphasizing the need for holistic evaluation of security decisions.

Recent work has identified security vulnerabilities in smart lock systems across several components. Ho et al. \cite{ho2016smartlocks} analyze five commercial smart locks and show that weaknesses can arise not only from implementation bugs, but also from system design, revocation handling, access logging, and automatic unlocking interaction models. Caballero-Gil et al. \cite{caballero2023smartlocks} and Turk \cite{turk2023smartlock} report weaknesses in Bluetooth communication, including replay and interception attacks, while Diao et al. \cite{diao2025masterlock} identify vulnerabilities in device protocols that allow unauthorized access persistence and manipulation of access rights. Veijalainen and Karlsson \cite{veijalainen2021evaluating} complement these works with a threat-model-driven penetration test of a smart door lock system, deriving concrete tests from reconnaissance and STRIDE/DREAD-based threat modeling to analyze app-server communication, replay behavior, gateway exposure, and key-tag access. Together, these works show that vulnerabilities in smart lock systems can already arise at the device and communication layer, before any cloud-side access control logic is even involved.

Beyond the device and communication layer, vulnerabilities have also been identified in the companion mobile applications that control smart locks. Allen et al. \cite{allen2024androidiot} highlight security flaws in companion mobile applications, such as exposure of API endpoints and authentication mechanisms. Complementing these findings, Wang et al. \cite{wang2019mirror} propose a scalable methodology to detect such issues through the analysis of mobile companion applications, showing that insecure design patterns can propagate across multiple devices due to code reuse. Together with the device-level findings above, these results demonstrate that vulnerabilities can occur at different layers of smart lock systems, motivating the need for stronger access control mechanisms.

User-oriented research further shows that smart lock security cannot be evaluated only from a technical perspective. Hazazi and Shehab \cite{hazazi2024automation} investigate smart lock automation in smart home environments and show that users perceive such automation in terms of security, convenience, awareness, privacy, and reliability. This is relevant for access-control systems because stronger verification mechanisms must not only improve protection, but also remain understandable and acceptable for users.

% Recent work has identified security vulnerabilities in smart lock systems across several components. Caballero-Gil et al. \cite{caballero2023smartlocks} and Turk \cite{turk2023smartlock} report weaknesses in Bluetooth communication, including replay and interception attacks, while Diao et al. \cite{diao2025masterlock} identify vulnerabilities in device protocols that allow unauthorized access persistence and manipulation of access rights. In addition, Allen et al. \cite{allen2024androidiot} highlight security flaws in companion mobile applications, such as exposure of API endpoints and authentication mechanisms. Complementing these findings, Wang et al. \cite{wang2019mirror} propose a scalable methodology to detect such issues through the analysis of mobile companion applications, showing that insecure design patterns can propagate across multiple devices due to code reuse. Together, these works demonstrate that vulnerabilities can occur at different layers of smart lock systems, motivating the need for stronger access control mechanisms.

Two complementary implementation-oriented studies investigate the integration of smart lock systems with cloud infrastructures. Salvadore \cite{salvadore2022smartlocks} proposes a cloud-based orchestration system built on API-based integration with smart lock platforms such as Nuki, where the Web API is used to generate and manage time-limited virtual keys based on booking data. This enables automated access provisioning and reduces manual intervention. However, this approach relies mainly on booking data and does not evaluate additional verification at access time.

In contrast, Zebic \cite{zebic2023azure} examines the integration of a Nuki smart lock into Microsoft Azure from a process-oriented perspective. The work identifies and evaluates the individual steps required for integration, such as configuring authentication parameters, establishing API connections, and managing device-specific settings. By introducing metrics to assess the complexity of these steps, the study shows that integrating smart locks into cloud environments involves multiple configuration tasks that go beyond simple API usage.

Other prototype-based studies illustrate how smart locks can be implemented using embedded devices and cloud services. Maurya et al. \cite{maurya2020cloudnative} implement a cloud-native door lock system using Raspberry Pi, Arduino, Alexa, and Amazon Web Services, including services for image storage, face recognition, messaging, and device communication. Derbali \cite{derbali2023iotsdl} proposes an Internet of Things smart door lock system with temporary PIN codes, Radio Frequency Identification, Bluetooth-based access, remote unlocking, and access records. Zheng et al. \cite{zheng2024stm32} present an STM32-based smart door lock with multiple unlocking methods, including RFID, fingerprint recognition, keypad access, remote unlocking, and a vibration alarm. These works demonstrate the technical feasibility of cloud-connected and embedded smart lock systems.

While existing work has extensively analyzed security vulnerabilities, usability concerns, and system integration challenges in smart lock environments, the infrastructure cost associated with stronger access control mechanisms has not been studied as extensively. Additional verification steps, such as contextual validation or biometric verification, can improve access control but also introduce more complex system interactions, additional cloud service usage, and increased resource consumption.

To address this gap, this position paper focuses on the cost and infrastructure impact of stronger access control mechanisms in a cloud-managed smart lock system. By comparing three access control scenarios with increasing security levels, ranging from direct keypad access to booking-code validation and face verification, this work evaluates how stronger access control affects operational cost, CPU utilization, and memory utilization in a realistic cloud-based booking system.